\subsection*{Problem 2.1}

\subsubsection*{(a)}

Both frameworks implement \textbf{Dynamic Epistemic Logic (DEL)}. 
They start with an epistemic model (a Kripke structure) 
representing the agents’ knowledge and then update it when events occur. 
The differences lie in the kinds of events they support and the 
underlying data structures.

\paragraph*{\texttt{demo\_s5} – Public announcements only}

\subparagraph*{Representation.}
\begin{itemize}
    \item \textbf{States:} any hashable objects (e.g., tuples, integers).
    \item \textbf{Agents:} identified by integers (e.g., $a=0$, $b=1$).
    \item \textbf{Valuation:} a mapping from states to a set 
    of propositional atoms. The logic is \textbf{2‑valued} 
    (True/False) – every proposition is either true or false at every world.
    \item \textbf{Accessibility:} stored as \textbf{partitions} 
    (equivalence relations). For each agent, a list of blocks 
    (each block is a list of states). 
    This enforces S5 properties (reflexive, transitive, symmetric).
    \item \textbf{Formulas:} \texttt{Top}, \texttt{Info(world)} 
    (true only at that world), \texttt{Ng} (negation), 
    \texttt{Conj}, \texttt{Disj}, \texttt{Kn(agent,form)} (knowledge).
\end{itemize}

\subparagraph*{Event handling.}
Only \textbf{public announcements} are supported via \\ 
the function \texttt{upd\_pa(model, form)}. It:
\begin{enumerate}
    \item Keeps only those states where 
    the announced formula is true.
    \item Restricts each agent’s partition to the remaining 
    states (drops worlds and updates blocks).
    \item Keeps only actual worlds that survive.
\end{enumerate}
No other actions (private messages, factual changes, etc.) 
can be expressed.

\subparagraph*{{Solving a DEL problem.}}
In \texttt{demo\_s5}, the initial epistemic model 
$M$ must be \textbf{explicitly constructed} to encode 
all relevant information of the puzzle. 
For Cheryl’s birthday, $M$ contains exactly the ten 
candidate dates, with Albert’s relation grouping 
by month and Bernard’s relation grouping by day. 
The actual world is chosen from among these states. 
There is no way to “set” the real situation later – 
everything is given into the initial model.
A DEL problem in this framework is specified by:
\begin{enumerate}
    \item An initial epistemic model $M$ (
        states, partitions for each agent, valuation, actual worlds).
    \item A sequence of public announcements $\phi_1, \phi_2, \dots, \phi_n$.
\end{enumerate}
The solution proceeds by repeatedly applying \texttt{upd\_pa}:
\[
M^{(0)} = M,\qquad M^{(i+1)} = \texttt{upd\_pa}(M^{(i)}, \phi_{i+1}).
\]
After each announcement, the model shrinks: 
worlds where the announced formula is false disappear. 
The final model $M_n$ contains exactly 
those worlds that are consistent with all announcements. 
If the puzzle asks “what is the real world?” (e.g., Cheryl’s birthday), 
the answer is the unique actual world in $M_n$ (if exactly one remains). 
If the puzzle asks “does agent $a$ know $\psi$?”, 
one checks $\texttt{is\_true}(M_n, \texttt{Kn}(a,\psi))$.

\hl{0.5}

\paragraph*{\texttt{demo\_light} – Action models with vocabulary change}

\subparagraph*{Representation.}
\begin{itemize}
    \item \textbf{States:} integers (after bisimulation) or any hashable type.
    \item \textbf{Agents:} same as above.
    \item \textbf{Vocabulary (\texttt{voc}):} an explicit list of propositional 
    atoms that ``exist''. Truth evaluation is \textbf{3‑valued} 
    (True / False / None) – a proposition not in \texttt{voc} is undefined. 
    This allows modelling ignorance about the existence of a fact.
    \item \textbf{Valuation:} mapping from states to the set of 
    atoms true there (only atoms from \texttt{voc} appear).
    \item \textbf{Accessibility:} stored as a \textbf{list of triples} 
    \texttt{(ag, state1, state2)} – \\ the states(\texttt{state1,state2}) 
    that agent(\texttt{ag}) can't 
    distinguish.
    \item \textbf{Formulas:} \texttt{Top}, 
    \texttt{Prp(atom)}, \texttt{Neg}, 
    \texttt{Conj}, \texttt{Disj}, \texttt{K(agent,form)}, 
    \texttt{CK(agents,form)} (common knowledge).
\end{itemize}

\subparagraph*{Event handling – Action Models with Vocabulary Change (ACM).}

An \textbf{ACM} is a generalisation of a public announcement. It consists of:
\begin{itemize}
    \item \texttt{states}: a set of event indices (e.g., $0,1,2$).
    \item \texttt{pre}: a precondition formula for each event.
    \item \texttt{subst}: for each event, a list of pairs 
    \texttt{(proposition, formula)}. When the event occurs, 
    each proposition is \textbf{replaced} by the given formula 
    (evaluated in the original world). This models \textbf{factual change} 
    (e.g., flipping a switch, painting a wall).
    \item \texttt{rels}: a list of triples \texttt{(agent, event1, event2)} 
    describing which events an agent cannot distinguish.
    \item \texttt{actual}: the event(s) that actually happen.
\end{itemize}
Product update \texttt{upc(m, facm)} (or \texttt{upd} 
followed by bisimulation) computes the new epistemic model:
\begin{enumerate}
    \item New states are pairs \texttt{(world, event)} 
    where the precondition of the event holds at the world.
    \item New vocabulary = old vocabulary $\cup$ all 
    propositions appearing in any substitution of the ACM.
    \item New valuation: for each \texttt{(w,s)}, a 
    proposition $p$ is true iff its substitution 
    (or $p$ itself if no substitution) evaluates to true 
    at $w$ ($m,w \models p \text{ or }sub(p)$).
    \item New relations: \texttt{(w1,s1)} and 
    \texttt{(w2,s2)} are related for agent $ag$ iff 
    $w1$ and $w2$ were related in $m$ and 
    $s1$ and $s2$ are related in the $facm$.
    \item Actual worlds: pairs of actual world and actual event.
\end{enumerate}

High‑level constructors (all return ACMs) include:
\begin{itemize}
    \item \texttt{public(form)} – single event, 
    precondition \texttt{form}, no substitution. (Public announcement)
    \item \texttt{group\_m(group, form)} – group 
    learns \texttt{form}; outsiders cannot distinguish it from \texttt{Top}.
    \item \texttt{message(agent, form)} – private message to one agent.
    \item \texttt{test(form)} – silent test 
    (no one observes the restriction).
    \item \texttt{info(group, form)} – 
    group learns \textbf{whether} \texttt{form} is true or false (two events).
    \item \texttt{p\_change(substitution)} – public 
    factual change (all agents see the substitution).
    \item \texttt{perception(agent, form, group)} – 
    agent and group observe the truth value of \texttt{form} (no change).
    \item \texttt{perceived\_change(agent, prp, form, group)} – 
    agent and group observe the old truth value of 
    \texttt{prp} and simultaneously change it to \texttt{form}.
\end{itemize}

\subparagraph{Solving a DEL problem.}
In \texttt{demo\_light}, the initial model does \textbf{not} 
need to be hand‑crafted for the specific puzzle. 
Instead, one typically starts from a 
\textbf{blissful‑ignorance model} generated 
by \texttt{initM(agents, props)}. 
This model contains all logically possible worlds 
(all truth assignments to the given propositions). 
The real situation is then imposed \textbf{dynamically} 
using ACM actions.For example:
\begin{itemize}
    \item \texttt{test(actual\_situation)} 
    restricts the model to worlds where the
    actual situation holds, but 
    no agent observes this restriction
     – it is a silent, unobserved factual setting.
    \item Further actions (\texttt{info}, 
    \texttt{message}, \texttt{group\_m}, 
    \texttt{p\_change}, \texttt{perception}, etc.) 
    refine the model by adding private observations,
    factual changes, or vocabulary modifications.
\end{itemize}
Thus, a DEL problem is specified by:
\begin{enumerate}
    \item An initial epistemic model $M$ 
    (states, relations, valuation, actual worlds).
    \item A sequence of actions $\alpha_1, \alpha_2, 
    \dots, \alpha_n$, where each $\alpha_i$ is an ACM 
    (or a high‑level constructor that produces one).
\end{enumerate}
The solution proceeds by repeatedly applying 
product update:
\[
M^{(0)} = M,\qquad M^{(i+1)} = \texttt{upd}(M^i, \alpha_{i+1})
\]
After each action, the model may grow 
(because of pairs with different events) 
or shrink (because of preconditions). 
The final model $M^{(n)}$ represents the 
epistemic situation after all actions. Queries 
(e.g., “does agent $a$ know $\psi$?”, 
“what is the value of proposition $p$?”) 
are answered by evaluating the relevant 
formula in $M^{(n)}$ using \texttt{is\_true}.

\subsubsection*{(b)}

The extra representational power of \texttt{demo\_light} 
comes from two key innovations:

\begin{enumerate}
    \item \textbf{Explicit vocabulary and 3‑valued logic.}
        \begin{itemize}
            \item \texttt{demo\_s5} uses a fixed set of propositions; 
            every proposition is either true or false at every world. 
            There is no notion of an ``undefined'' proposition.
            \item \texttt{demo\_light} maintains an explicit list 
            \texttt{voc} of atoms that exist. A proposition not 
            in \texttt{voc} evaluates to \texttt{None} (undefined). 
            This allows modelling situations where agents are not 
            even aware of a certain fact, or where new propositions 
            are introduced dynamically (e.g., through an ACM substitution).
        \end{itemize}
    
    \item \textbf{Action models with substitution (ACM).}
        \begin{itemize}
            \item \texttt{demo\_s5} can only perform 
            \textbf{public announcements} – a single event that all agents 
            know has occurred, and which only removes worlds.
            \item \texttt{demo\_light} supports \textbf{any action model}, including:
                \begin{itemize}
                    \item \textbf{Private and group announcements} 
                    (e.g., \texttt{message}, \texttt{group\_m}, \texttt{info}) 
                    – where different agents have different access to the events.
                    \item \textbf{Silent tests} (\texttt{test}) 
                    – the world changes but no agent observes it.
                    \item \textbf{Factual changes} (\texttt{p\_change}, 
                    \texttt{unobserved\_change}, \texttt{perceived\_change}) 
                    – the truth value of propositions can be altered, 
                    not just eliminated.
                    \item \textbf{Perception changes} – an agent 
                    observes the truth of a formula \\ 
                    (\texttt{perception}) or observes a factual change and \\ 
                    thereby learns the
                    old value (\texttt{perceived\_change}). \\
                    For example, if an agent sees a light switch being flipped, 
                    they learn that the light was previously in the opposite state.
                \end{itemize}
            \item Because public announcements are a special case of an ACM 
            (single event with no substitution), \texttt{demo\_light} 
            can simulate every scenario that \texttt{demo\_s5} can, 
            but the converse is false.
        \end{itemize}
\end{enumerate}

In summary, \texttt{demo\_light} is strictly more expressive: 
it can model public announcements, private messages, 
factual changes, unobserved changes, perception changes, 
and combinations thereof – all within a unified framework 
of action models with vocabulary change. \texttt{demo\_s5} 
is limited to public announcements only, making it suitable 
for puzzles like Cheryl’s birthday but inadequate 
for scenarios involving private information or world‑altering events.

\subsubsection*{(c)}

\paragraph*{(i) Can the Sally–Anne test be modeled using \texttt{demo\_s5}?}

\textbf{No.} The \texttt{demo\_s5} framework has two 
fundamental limitations that make it incapable of modeling the Sally–Anne scenario:

\begin{enumerate}
    \item \textbf{Only public announcements.}
    \texttt{demo\_s5} supports only one type of event:
    a truthful public announcement that all agents hear and 
    know that everyone hears. In the Sally–Anne test, Anne moves the 
    toy \textbf{privately} while Sally is out of the room. Sally 
    does not observe the move, so her belief remains outdated. 
    This is a private, unobserved action. Public announcements 
    cannot capture such asymmetric information updates.
    
    \item \textbf{No factual change.} 
    In \texttt{demo\_s5}, the valuation 
    (which propositions are true at each world) is fixed at 
    the start and never changes. The only operation is world elimination. 
    The Sally–Anne test requires a \textbf{factual change}: 
    the toy’s location changes from Sally’s basket to Anne’s box. 
    This cannot be represented by simply removing worlds; 
    the truth value of the proposition “toy in Sally’s basket” 
    must become false, and “toy in Anne’s box” must become true. 
    \texttt{demo\_s5} has no mechanism for substitution or valuation update.
\end{enumerate}

One might attempt to encode both possible locations in the initial 
model and then use a public announcement that “Anne moves the toy”. 
However, that announcement would be heard by Sally, making her 
aware of the move and destroying the false belief that is essential 
to the test. Therefore, \texttt{demo\_s5} cannot model the Sally–Anne test.

\paragraph*{(ii) Can it be modeled using \texttt{demo\_light}?}
\textbf{Yes.}

\texttt{demo\_light} can model the test using action 
models with vocabulary change (ACM). 
The implementation follows these steps (see \texttt{sally\_anne.py}):

\begin{enumerate}
    \item \textbf{Agents and propositions:} Sally ($a$) and Anne ($b$); 
        $P(1)$ = “toy in Sally’s basket”, $P(2)$ = “toy in Anne’s box”.
    \item \textbf{Initial model:} Start from a blissful‑ignorance model \\
        (\texttt{initM([sally, anne], [P(1), P(2)])}) 
        containing all $2^2=4$ worlds.
        Silently set the actual world (toy in Sally’s basket) using \\
        \texttt{test(Conj([in\_sally, Neg(in\_anne)]))}.
    \item \textbf{Private move action (ACM):} 
        Construct an action model with
        two events: $0$ (Anne moves the toy) and $1$ (no move).
        \begin{itemize}
            \item Preconditions: event $0$ requires \texttt{in\_sally}; 
                event $1$ requires $\top$.
            \item Substitutions: event $0$ maps $P(1)$ to $\texttt{False}$ 
                and $P(2)$ to $\texttt{True}$; event $1$ does nothing.
            \item Accessibility: Anne distinguishes $0$ from $1$; 
                Sally cannot distinguish them (edges $(0,1)$ and $(1,0)$).
            \item Actual event: $0$ (Anne moves the toy).
        \end{itemize}
        The ACM is applied with \texttt{upd(m1, private\_move())}.
    \item \textbf{Verification:} After the update, evaluate:
        \begin{itemize}
            \item $K(\text{Sally}, P(2))$ – 
            Sally does know the toy is in Anne’s box. (False)
            \item $K(\text{Anne}, P(2))$ – Anne knows the true location. (True)
            \item Sally considers it possible that the toy is still in her basket\\ 
                (i.e., $\neg K(\text{Sally}, \neg P(1))$), 
                capturing her false belief. (True)
        \end{itemize}
\end{enumerate}

The program outputs the expected truth values, 
confirming that \texttt{demo\_light} successfully models the false‑belief scenario.